\documentclass[11pt]{ctexart}
\usepackage[margin=2.3cm]{geometry}
\usepackage{booktabs}
\usepackage{array}
\usepackage{xcolor}
\usepackage{enumitem}
\usepackage{hyperref}
\hypersetup{colorlinks=true, linkcolor=blue!55!black, urlcolor=blue!55!black}
\setlist{itemsep=1pt, topsep=2pt}

\definecolor{cue}{RGB}{180,70,40}
\newcommand{\cue}[1]{{\color{cue}\textbf{【#1】}}}

\title{\textbf{汇报演讲稿}\\[3pt]
\large ETF--股指期货 期现基差套利复现}
\author{许哲圣 \quad 国泰海通自营}
\date{2026 年 6 月}

\begin{document}
\maketitle

\noindent\small\textit{用法：左手 report.pdf，右手这份讲稿。每段开头 \cue{红字} 是``翻到哪、指什么''。
全程约 9--10 分钟。加粗的是必须说出口的关键词。节奏：方法 2 分钟 $+$ 结果 3 分钟 $+$ 融券发现
1.5 分钟，是重头，别赶。}

\section*{0\quad 开场（30 秒）\quad \cue{封面 / 摘要}}
各位老师好，我汇报的课题是 \textbf{ETF 高频套利}，具体做的是\textbf{股指期货--ETF 的期现基差套利复现}。
我复现的是三篇卖方研报——\textbf{东证、银河、华泰柏瑞}——它们讲的是同一个框架。数据全部用
\textbf{tushare 真实 A 股数据}：上证50、沪深300、中证500、中证1000 四套股指期货加对应宽基 ETF，
区间 2019 到 2026。先把结论给出来：

\begin{quote}
\textbf{期现套利赚的是「年化基差 $+$ 分红 $-$ 成本」，而融券约束决定 A 股能不能真正执行。}
\end{quote}
这句是整个报告的主线，后面所有结果都在论证它。

\section*{1\quad 研究背景与目标（1 分钟）\quad \cue{第 1 节，表 1}}
为什么不能直接抄研报数字？两个坑：

第一，\textbf{分红基差}。高股息品种比如上证50、沪深300，不做分红调整，原始 $F-S$ 会
\textbf{系统性高估贴水}，信号失真。

第二，\textbf{融券约束}。A 股股指期货 2019--2026 长期贴水，要赚钱往往得做反向套利、也就是
\textbf{做空 ETF}，受融券券源制约。研报里漂亮的正收益，很大程度依赖融券资源。

\cue{指表 1} 这张表是我和原研报的口径对照——样本区间、现货端、记账方式、融券假设，说明哪些
和研报一致、哪些我有意做得更严格。

这里\textbf{诚实前置一个工具约束}：任务原本要求用 hftbacktest 和 Hummingbot，但二者和 A 股
ETF 套利有结构性错配——\textbf{Hummingbot 只能接加密交易所，连不了上交所深交所}；
\textbf{hftbacktest 需要逐笔 L2，而免费可回溯的 A 股 L2 实际不存在}。所以唯一能用真实 A 股数据
复现研报数值的就是纯实证回测，这就是主体 Track 0；hftbacktest 和 Hummingbot 我做成了执行
验证层，就是后面的 Track A 和 Track B。

\section*{2\quad 数据与方法（2 分钟，核心）\quad \cue{第 2 节，表 2 + 公式}}
\textbf{数据}\cue{指表 2}：全部来自 tushare，缓存成本地 parquet，可离线复跑。包括指数日线、
股指期货主力连续、ETF 日线复权、以及\textbf{全收益指数}——它和价格指数的差就是我估股息率的依据。
表 2 是四个品种对和年化股息率：上证50 是 3.4\%、沪深300 是 2.5\%、中证500 是 1.6\%、
中证1000 是 1.1\%。注意这个梯度，后面解释分红效应要用。

\textbf{基差模型}\cue{指公式 1、2}：两个公式。持有成本定价——$F$ 理论值 $= S\,(1+(r_f-d)\,t/365)$；
年化基差率 $=\frac{F-S}{S}\cdot\frac{365}{dte}$。关键在\textbf{分红调整}：持有``多 ETF、空期货''
到交割会额外收到成分股分红，而期货不付息，所以真正可交易的 carry 是原始基差\textbf{再加股息率}。
这修正了高股息品种高估贴水的偏差。

\textbf{交易信号}：三种信号都是\textbf{因果、滚动、无前视}，各对应一篇研报：
\begin{itemize}
  \item \textbf{conv 基差收敛}（银河/东证）：先识别升水/贴水/中性 regime，分红调整年化基差超
        $+3\%$ 入场，持有整段升水 regime，跌破 $-0.5\%$ 才平。低换手，对应研报约 40 天持仓。
  \item \textbf{galaxy 分位数}（银河）：年化基差率滚动历史分位，超 80 分位做正向。
  \item \textbf{orient 均值回归}（东证）：年化基差率 z-score。
\end{itemize}

\textbf{收益与成本}：期现套利\textbf{到期锁定终值}，一笔在某年化利率入场的交易持有期内近似稳定
赚这个利率——这就是研报``胜率近 100\%、回撤极小''的来源。我据此\textbf{入场锁定 carry} 逐日累计。
还要强调\textbf{分红不重复计}：carry 用不含分红的原始基差；分红通过 ETF 全收益腿、按 A 股实际
\textbf{分红季节 6--7 月}单独捕获；分红调整基差只用于入场判断，不重复进盈亏。

\section*{3\quad 复现结果（3 分钟，重头戏）}

\textbf{3.1 主结果}\cue{第 3 节，表 3 + 图 1、图 2}。先看融券开启情景，即复现研报、假设机构有融券
资源。四品种年化收益\textbf{全部进入研报区间}，看最后一行——\textbf{多品种复合：年化 5.6\%、夏普
1.94、最大回撤 $-3.5\%$}。银河研报是 6.6\% / 夏普 1.78 / $-3.4\%$，\textbf{高度吻合}，框架复现成功。
\cue{翻图 1} 复合净值：蓝线融券开启稳健上行到 1.44；红线融券受限一路下行到 0.90——先记住这个对比，
第 4 节回来。\cue{翻图 2} 四品种各自净值，回撤都在 5\% 以内。

\textbf{3.2 分红效应}\cue{翻图 3}。这张图是分红调整的直接证据：沪深300 年化基差率，原始口径看着
是``贴水'',分红调整后系统性上移约 2.5\%——正好是它的股息率，把可交易 carry 还原出来了。

\textbf{3.3 动态 ETF 选择}\cue{表 5}。东证精修：每个指数多只候选 ETF，每月按``流动性减跟踪误差''
择优。提升最明显的是 \textbf{IM 小盘}（候选 ETF 分散度最大），年化 7.8\%$\to$8.2\%、夏普
1.74$\to$1.85；IF 同质化高几乎没变；IH 只有一只 ETF 完全不变。和东证``动态选择更稳定''一致。

\textbf{3.4 逐合约严格记账——诚实加分项}\cue{表 6 + 图 4}。前面 5.6\% 是``连续主力 $+$ 锁定 carry''
近似。为检验它，我又做了\textbf{逐合约持有至交割}严格记账：单一合约内逐日盯市持有到交割，
\textbf{无换月跳价、终值在交割真实收敛}。\cue{指表 6} 严格记账复合年化 \textbf{3.9\%}，比近似的
5.6\% \textbf{更保守}，说明锁定 carry 近似略偏乐观。\cue{翻图 4} 红线（近似）始终略高于蓝线（严格），
但走势一致——\textbf{验证近似可用但偏乐观}。我想强调：\textbf{我没有只报漂亮的 5.6\%，而是用两套
口径互相印证}——宽松口径对上研报，严格口径给出更保守、更接近 4\% 的真实可执行收益。

\section*{4\quad 关键发现：融券敏感性（1.5 分钟，最有洞察）\quad \cue{表 7}}
\cue{指表 7} 回到开头那句话的后半句。我把反向套利仓位按融券可得比例缩放、再收融券费做敏感性：
\begin{itemize}
  \item \textbf{完全没有融券：复合年化只有 0.2\%，夏普 $-1.3$。}
  \item 完全可融券、零费：5.6\%。
  \item 完全可融券、但融券费 6\%：从 5.6\% 掉到 \textbf{2.8\%}。
\end{itemize}
\cue{回指图 1 红线} 这就解释了那条往下走的红线。\textbf{结论}：研报收益不是单纯的基差信号收益，
而是\textbf{高度依赖融券资源}。对自营台的现实含义——这策略能不能做，取决于\textbf{能不能借到券}，
而不只是基差信号好不好看。正好印证华泰柏瑞``融券约束制约反向套利''。

\section*{5\quad 两个执行验证层（1.5 分钟）}
\textbf{Track A：执行撮合}\cue{表 8 + 图 5}。hftbacktest 风格的执行 proxy：每次调仓拆 8 个子单，
模拟被动排队、超时吃价、双腿滑点。执行成本逐日回灌后，复合年化从 5.6\% 降到 \textbf{4.6\%}、夏普
1.94 降到 1.39。\textbf{执行成本拖累可量化}。盘口已抽象成可插拔接口，ETF 五档现在用 tushare 已能录，
股指期货完整五档等券商 CTP。

\textbf{Track B：币安现货--永续}\cue{表 9 + 图 6}。期现套利在加密市场的等价物——\textbf{现货 vs
永续}的 cash-and-carry，靠\textbf{资金费}吃 carry，对应 A 股``多 ETF、空期货''。数据全是币安免费
公共行情，复用同一套信号。复合年化 7.1\%、夏普 5.24。资金费 carry 的\textbf{低波动高夏普}和 A 股
正向套利一致——跨市场相互印证。配套给了 Hummingbot 纸面部署制品。

\section*{6\quad 关键发现与局限（45 秒）\quad \cue{第 4 节}}
四条，前两条是发现、后两条是主动说的局限：
\begin{enumerate}
  \item \textbf{融券约束是 A 股期现套利的核心制约}——最重要的现实结论。
  \item \textbf{分红基差不可忽略}——高股息品种不调整就失真。
  \item 和研报差异来自样本区间、ETF 选择、口径，不是缺陷，复合层面吻合度最高。
  \item 建模近似：主力连续、锁定 carry 假设收敛单调、股息率用全收益指数估计。
\end{enumerate}

\section*{7\quad 后续计划 $+$ 收尾（30 秒）\quad \cue{第 5 节}}
后续三条线：Track 0 把 regime 纳入月度监控；Track A 接券商 CTP 五档做真实事件回放；Track B 在
paper/testnet 实跑挂机。收尾：
\begin{quote}
\textbf{这个项目我不只是复现了研报数字，更重要的是搞清楚了背后两个前提——分红要调整、融券是命门。
复合层面我对上了银河的 6.6\%，同时用严格记账诚实地给出了更保守的 3.9\%。} 谢谢老师，请批评指正。
\end{quote}

\section*{附：被追问时的标准答案}
\small
\begin{tabular}{>{\raggedright\arraybackslash}p{4.6cm} >{\raggedright\arraybackslash}p{10.5cm}}
\toprule
老师可能问 & 你这样答 \\
\midrule
单品种数字对不齐？（如 IM 报告 1.9\% 你 7.8\%） & 样本区间不同（研报到 2024.4，我到 2026）、ETF 选择和口径不同；\textbf{复合层面是吻合的}，单品种差异在预期内，第 4 节第 3 条专门说了。 \\
分红会不会重复计？ & 不会。carry 用\textbf{原始基差}，分红单独通过 ETF 全收益腿按 6--7 月实现季节捕获，分红调整基差只用于入场信号。 \\
有没有前视偏差？ & 没有。所有信号滚动因果，分位数和 z-score 只用历史窗口，min\_periods 控制。 \\
为什么逐合约更可信？ & 单一合约内盯市持有到交割，\textbf{无换月跳价}，终值在交割真实收敛，回撤更真实。 \\
真实 L2 数据有吗？ & ETF 腿现在用 tushare 五档已能录制；股指期货完整五档需券商 CTP，接口已抽象就绪。 \\
为什么做 Track B 加密的？ & 任务点名要 Hummingbot，但它连不了 A 股；现货--永续是期现套利最干净的加密类比，能在免费数据上跑出真实净值、跨市场印证资金费 carry 特征。 \\
代码可靠吗？ & 41 个单元测试全通过，config.py 是单一可信源，回测和报告读同一套参数。 \\
\bottomrule
\end{tabular}

\end{document}
